% !TEX program = lualatex
\documentclass[12pt, a4paper]{article}
\usepackage{master}

\DeclareWorkGR{Cat}{범주론}{Κατηγορίαι}{Categoriae}
\DeclareWorkGR{Top}{토피카}{Τοπικά}{Topica}
\DeclareWorkGR{APo}{분석론 후서}{Ἀναλυτικὰ ὕστερα}{Analytica posteriora}

\fancyhead[L]{\footnotesize\color{midgray}오르가논 강독}
\fancyhead[R]{\footnotesize\color{midgray}제7주}

\begin{document}
\thispagestyle{firstpage}

\weeksection{제7주}{이사고게 IV--V장}{종차와 고유성}

\subsection{종차: 종을 만드는 것 (IV장)}

\subsubsection{종차의 세 가지 의미}

종차(\gr{διαφορά})는 이사고게에서 철학적으로 가장 복잡한 장입니다. 포르피리오스는 종차를 세 가지 의미로 구별합니다.

첫째, 일반적(\gr{κοινῶς}) 의미의 종차입니다. 한 사물이 다른 사물과 구별되는 어떤 특징이든 종차가 될 수 있습니다 --- 분리 가능한 부수성까지 포함합니다. 소크라테스는 지금 앉아 있다는 점에서 플라톤과 다를 수 있습니다. 이 의미의 종차는 논리학에서는 쓸모가 없습니다 --- 너무 넓기 때문입니다.

둘째, 고유한(\gr{ἰδίως}) 의미의 종차입니다. 분리 불가능한 특징으로 구별하는 것입니다 --- 예컨대 매부리코인 것(\gr{τὸ γρυπόν}). 이것은 소크라테스를 다른 사람들과 구별하지만, 사람이라는 종의 본질에 기여하지는 않습니다.

셋째, 가장 고유한(\gr{ἰδιαίτατα}) 의미의 종차입니다. 이것이 기술적 의미의 종차이며, 종을 그것이 무엇이게 만드는 종차입니다. 예컨대 `이성적인'(\gr{λογικόν})은 사람이라는 종의 가장 고유한 종차입니다 --- `동물'(류) + `이성적인'(종차) = `이성적 동물'(사람의 정의)이기 때문입니다.

\subsubsection{종차의 두 가지 역할: 분할하는 것과 구성하는 것}

포르피리오스는 종차에 두 가지 역할을 부여합니다. 한편으로 종차는 류를 분할합니다(divisivae) --- `이성적인'과 `비이성적인'은 `동물'이라는 류를 두 하위 부류로 나눕니다. 다른 한편으로 종차는 종을 구성합니다(constitutivae) --- `이성적인'은 `동물'과 결합하여 `사람'이라는 종을 만들어냅니다. 분할적 종차와 구성적 종차는 같은 종차를 서로 다른 방향에서 본 것입니다.

6주차에서 다룬 포르피리오스의 나무에서 이 이중 역할이 직관적으로 이해됩니다. 나무를 위에서 아래로 읽으면, 각 분기점에서 종차가 류를 분할합니다. 나무를 아래에서 위로 읽으면, 종차들이 축적되어 최저종의 정의를 구성합니다 --- 사람 = 실체이면서 + 물체적이면서 + 생명 있으면서 + 감각 있으면서 + 이성적인 것.

\subsubsection{종차는 어떤 범주에 속하는가}

종차의 범주적 지위 문제는 아리스토텔레스 해석자들 사이에서 오래된 난제입니다. 포르피리오스는 종차가 `어떤 종류의 것인지'(\gr{ἐν τῷ ὁποῖόν τί ἐστι})를 말해준다고 합니다. 이것은 류의 `무엇인지'(\gr{ἐν τῷ τί ἐστι})와 대조됩니다. `소크라테스는 무엇인가?' --- `동물'(류). `소크라테스는 어떤 종류의 동물인가?' --- `이성적인'(종차).

그런데 `어떤 종류의 것'이라는 표현은 질(\gr{ποιόν})을 연상시킵니다. 실제로 아리스토텔레스는 \citework{Cat}{3b10 이하}에서 종차를 질의 사례로 제시하는 것처럼 보입니다.\footnote{Ackrill, \gri{Aristotle's Categories and De Interpretatione}, 3b10 이하 주석.} 만약 종차가 질이라면, 범주론의 체계에서 종차는 `주어 안에 있는 것'(부수성)이어야 합니다. 그런데 종차는 본질적입니다 --- `이성적인'은 사람이 사람이기 위한 조건이지, 사람에게 부수적으로 붙은 속성이 아닙니다.

보에티우스(Boethius)는 이 문제를 `실체 안의 질'(quale in substantia)과 `실체 밖의 질'(quale non in substantia)의 구별로 해결하려 했습니다. 그러나 이 구별은 포르피리오스의 텍스트에 명시적 근거가 없으며, 범주론의 사중 분류에도 깔끔하게 맞지 않습니다. 아비첸나(Ibn S\={\i}n\=a)는 나중에 포르피리오스의 종차 정의 자체를 정면으로 비판합니다.\footnote{Cristina Cerami, ``Avicenna against Porphyry's Definition of Differentia Specifica,'' \gri{Documenti e studi sulla tradizione filosofica medievale} 26 (2015): 129--183.}

여기서 사고 연습을 해봅시다. 2주차의 범주론 사중 분류를 떠올려 봅시다. `주어에 대해 말해지면서 주어 안에 있지 않은 것'(제2실체), `주어 안에 있으면서 주어에 대해 말해지지 않는 것'(개별적 속성). 종차는 어디에 놓아야 할까요? 종차는 분명 주어에 대해 말해집니다(사람은 이성적이다) --- 그러나 `이성적인'은 제2실체가 아닙니다. 종차는 주어 안에 있다고 말하기도 어렵습니다 --- 그것은 부수성에 해당하니까요. 범주론의 사중 분류는 종차를 위한 자리를 명확히 마련하지 않았으며, 이것이 이사고게가 범주론에 `추가'해야 했던 이유 가운데 하나입니다.

\subsubsection{류 + 종차 = 정의}

류 + 종차 = 정의라는 공식이 작동하려면 두 가지 조건이 필요합니다. 첫째, 종차는 류보다 좁아야 합니다 --- `이성적인'은 `동물' 전체가 아니라 동물의 일부에만 적용됩니다. 둘째, 종차는 류와 다른 차원에서 와야 합니다 --- 류(`동물')가 `무엇인지'를 말해주면, 종차(`이성적인')는 `어떤 종류인지'를 말해줍니다. 류 안에서 또 다른 류를 종차로 사용하면 정의가 순환합니다.

이 공식은 나중에 \citework{APo}{II권}에서 정의와 논증의 관계를 분석할 때 핵심적인 역할을 합니다. 분석론 후서에서 아리스토텔레스는 정의가 논증의 `원리'(\gr{ἀρχή})라고 주장합니다 --- 정의에서 출발하여 논증이 구성됩니다. 포르피리오스의 이사고게는 이 정의의 구성 요소(류 + 종차)를 체계적으로 분석함으로써, 분석론 후서를 위한 예비 작업이기도 합니다.

\subsection{고유성: 본질을 따르되 본질이 아닌 것 (V장)}

\subsubsection{고유성의 네 가지 의미}

포르피리오스는 아리스토텔레스의 \citework{Top}{I.5, 102a18--30}를 따라 고유성(\gr{ἴδιον})의 네 가지 의미를 구별합니다.

{\footnotesize
\begin{concepttable}{|c|c|c|c|l|}
\hline
\rowcolor{tableheadblue}
{\color{h1navy}\bfseries 의미} &
{\color{h1navy}\bfseries 해당 종에만?} &
{\color{h1navy}\bfseries 모든 구성원에?} &
{\color{h1navy}\bfseries 항상?} &
{\color{h1navy}\bfseries 예시} \\ \hline
첫째 & 예 & 아니오 & — & 의술 \\ \hline
둘째 & 아니오 & 예 & — & 두 발 달림 \\ \hline
셋째 & 예 & 예 & 아니오 & 노년의 백발 \\ \hline
넷째 & 예 & 예 & 예 & 웃을 수 있음 (\gr{τὸ γελαστικόν}) \\ \hline
\end{concepttable}
}

넷째 의미만이 기술적 의미의 고유성입니다 --- 세 가지 조건(`모든 것에', `그것에만', `항상')을 모두 만족해야 합니다.

\subsubsection{대표적 예시: 웃을 수 있음(\gr{τὸ γελαστικόν})}

모든 사람은 웃을 수 있고, 사람만이 웃을 수 있으며, 사람은 항상(원리적으로) 웃을 수 있습니다. `항상'의 의미를 정확히 이해해야 합니다 --- 지금 실제로 웃고 있지 않더라도, 웃을 수 있는 능력은 사람에게 항상 속합니다. 이것은 현실태(\gr{ἐνέργεια})가 아니라 가능태(\gr{δύναμις})의 차원에서 `항상'입니다. 4주차에서 다루었던 질 범주의 첫 번째 종류, 즉 상태(\gr{ἕξις})와의 연결이 여기서 분명해집니다 --- 웃을 수 있음은 사람이라는 종에 속하는 안정적 역량입니다.

\subsubsection{고유성과 정의의 차이}

넷째 의미의 고유성은 종과 역전됩니다(\gr{ἀντιστρέφει}) --- `사람이면 웃을 수 있고, 웃을 수 있으면 사람이다.' 외연이 동일합니다. 정의도 종과 외연이 동일합니다. 그렇다면 고유성과 정의는 어떻게 다를까요?

차이는 이것입니다: 정의는 종의 본질을 말해주지만(\gr{τὸ τί ἦν εἶναι}), 고유성은 본질을 말해주지 않습니다. `사람은 이성적 동물이다'는 사람이 \emph{왜} 사람인지를 설명합니다. `사람은 웃을 수 있다'는 사람에게 항상 따라다니는 것이 \emph{무엇인지}를 말해주지만, 사람이 왜 사람인지를 설명하지 않습니다. 스콜라 철학의 용어를 빌리면, 종차는 본질의 구성적 요소(essentiale constituens)이고, 고유성은 본질의 결과적 요소(essentiale consequens)입니다.\footnote{\gri{New Catholic Encyclopedia}, ``Property (Logic).''}

사람이 웃을 수 있는 것은 사람이 이성적이기 때문입니다 --- 이성이 있기에 유머를 이해하고 웃을 수 있습니다. 고유성은 정의(본질)로부터 파생됩니다. 그러나 `웃을 수 있음'이라는 사실만 가지고는 사람의 본질이 무엇인지를 역추적할 수 없습니다 --- 정의는 원인을 말해주고, 고유성은 결과만 말해줍니다.

여기서 사고 연습을 해봅시다. 현대 생물학에서 `인간의 DNA 서열'은 고유성일까요, 종차일까요? 네 번째 의미의 고유성 조건을 만족합니다. 그러나 DNA 서열이 인간의 본질을 \emph{구성한다}고 말할 수 있을까요? 아리스토텔레스적 틀에서는, 이성(\gr{λόγος})이 인간의 본질적 종차이고 DNA는 그 본질의 물질적 기반(결과적 요소)일 것입니다.

% ── 참고문헌 ──
\subsection*{참고문헌}
\begin{references}

\refitem Ackrill, J.\ L. \gri{Aristotle's Categories and De Interpretatione}. Clarendon Aristotle Series. Oxford: Clarendon Press, 1963.

\refitem Barnes, Jonathan. \gri{Porphyry: Introduction}. Clarendon Later Ancient Philosophers. Oxford: Clarendon Press, 2003.

\refitem Busse, Adolf, ed. \gri{Porphyrii Isagoge et in Aristotelis Categorias Commentarium}. CAG IV.1. Berlin: G.\ Reimer, 1887. Eng.\ trans.: Jonathan Barnes, \gri{Porphyry: Introduction}, Clarendon Later Ancient Philosophers. Oxford: Clarendon Press, 2003.

\refitem Cerami, Cristina. ``Avicenna against Porphyry's Definition of Differentia Specifica.'' \gri{Documenti e studi sulla tradizione filosofica medievale} 26 (2015): 129--183.

\end{references}

\end{document}
